% Based on templates/tutorial.tex; finished notes omit dialogue and check records.
% Source: LEC07 Imaging Geometry - Tutorial Problems.pdf, PDF p. 1, Q1--Q3.
% Supplement: LEC07 Imaging Geometry - 2026.pdf, pp. 10--12, 26--32.
% Study date: 2026-11-03, Week 12; this does not establish an official timetable change.
% No tutorial answer key supplied. Proofs and degeneracy conditions derived here.

\begingroup
\setlength{\parskip}{0.4em}
\exercisegroup{SC4061-TUT07}{Tutorial 07：Imaging Geometry}
\label{tutorial:SC4061-TUT07}

\noteintro{
{\small 学习日期：2026-11-03（第十二教学周）。来源：\emph{LEC07 Imaging Geometry -- Tutorial Problems}，PDF p.~1，Questions 1--3；标定流程参照主讲义 pp.~26--32。原题未附参考答案；以下给出几何证明、参数估计流程及适用条件。}
}

\exercisequestion{SC4061-TUT07-Q01}{共线空间点的投影仍共线}

\textbf{对应讲义：}\relatedtopic{SC4061-L11-T02}{针孔模型与透视投影}。

\textbf{题目。}证明：在 pinhole camera（针孔相机）中，三个共线的空间点投影后仍然共线。

\begin{checkedsolution}
采用理想针孔模型，设成像平面为 $\Sigma$，光心为 $O$，三个空间点 $P_1,P_2,P_3$ 位于直线 $L$ 上，且均有有限的像点。先考虑 $O\notin L$ 的一般情况。

\begin{enumerate}[label=\arabic*.]
  \item 直线 $L$ 与直线外的光心 $O$ 唯一确定一个平面 $\Pi$。
  \item 投影线 $OP_1,OP_2,OP_3$ 均位于 $\Pi$ 内。像点 $p_i=OP_i\cap\Sigma$ 因而同时属于 $\Pi$ 和 $\Sigma$。
  \item $\Pi$ 包含光心，而成像平面 $\Sigma$ 不包含光心，所以两平面不同；已有有限像点保证它们相交，其交集为直线 $\ell$。于是 $\boxed{p_1,p_2,p_3\in\ell=\Pi\cap\Sigma}$，即三个像点共线。
\end{enumerate}

\textbf{退化情况。}若 $L$ 经过光心，所有能成像的点重合在同一像点。镜头畸变可能使实际照片中的直线弯曲，本题结论针对理想针孔投影。
\end{checkedsolution}

\exercisequestion{SC4061-TUT07-Q02}{同方向平行线与消失点}

\textbf{对应讲义：}\relatedtopic{SC4061-L11-T02}{透视投影与深度}。

\textbf{题目。}走廊照片中，沿纵深延伸的扶手、墙边和地板接缝在空间中相互平行，图像中的延长线却汇聚到 vanishing point（消失点）。解释原因。

\begin{checkedsolution}
\nopagebreak[4]
\textbf{几何解释。}这组直线位置不同，但具有共同方向 $\bm d$。沿它们向远处观察时，空间点与光心的连线趋向经过光心、方向为 $\bm d$ 的同一条直线。该直线与成像平面的交点就是共同的消失点。固定相机后，消失点由这组直线的方向决定，与各条直线的位置无关。

\begin{center}
  % Original schematic of Tutorial 7 Q2, not a reproduction of the corridor photo.
% Solid segments are finite visible portions; dashed lines extend to V.
\begin{tikzpicture}[x=1cm,y=1cm,font=\small,>=Latex]
  \node[anchor=south] at (1.6,2.55) {空间中：共同方向 $\bm d$};
  \foreach \x in {0.4,1.6,2.8}{
    \draw[noteBlue,thick,->] (\x,0) -- (\x,2.3);
  }
  \node[anchor=north,font=\footnotesize] at (1.6,-0.12) {直线相互平行，位置不同};

  \draw[->,thick,black!60] (3.5,1.15) -- (4.7,1.15)
    node[midway,above,font=\footnotesize] {投影};

  \begin{scope}[shift={(5.4,0)}]
    \node[anchor=south] at (2,2.55) {图像中：共同消失点};
    \coordinate (V) at (2,2.3);
    \foreach \x in {0.3,2,3.7}{
      \coordinate (A) at (\x,0);
      \coordinate (B) at ($(A)!0.72!(V)$);
      \draw[noteBlue,thick] (A) -- (B);
      \draw[noteBlue,dashed] (B) -- (V);
    }
    \fill[ntuRed] (V) circle (2pt);
    \node[ntuRed,anchor=west] at (2.13,2.3) {$V$};
    \node[anchor=north,font=\footnotesize] at (2,-0.12) {像线延长后汇聚于同一点};
  \end{scope}
\end{tikzpicture}

  \par\smallskip
  \refstepcounter{figure}\label{fig:sc4061-tut07-vanishing}
  {\small 图~\thefigure：按原题走廊情境原创绘制；右图虚线为像线延长部分，$V$ 为共同消失点。}
\end{center}

\textbf{用投影公式核对。}在相机坐标系中，将任一条直线写为 $\bm P(\tau)=(X_0,Y_0,Z_0)^T+\tau(d_x,d_y,d_z)^T$。
取向远处延伸的方向，使 $d_z>0$。在主点为原点的图像物理坐标中，焦距为 $f>0$，且所考虑点的深度 $Z_0+\tau d_z>0$：
\[
  x(\tau)=f\frac{X_0+\tau d_x}{Z_0+\tau d_z},\qquad
  y(\tau)=f\frac{Y_0+\tau d_y}{Z_0+\tau d_z}.
\]
将分子、分母同除以 $\tau$ 并令 $\tau\to\infty$，得到
\[
  \boxed{(x_\infty,y_\infty)
    =\left(f\frac{d_x}{d_z},\ f\frac{d_y}{d_z}\right)}.
\]
极限只含方向分量与焦距，因此同方向的直线有相同的消失点。这里的 $(x,y)$ 是图像物理坐标；转为像素坐标还需处理像素尺寸、轴方向和主点偏移。

\textbf{方向与边界条件。}不同方向可以对应不同消失点。若 $d_z=0$，这组直线平行于成像平面，其可见投影仍平行，没有有限交点，称为消失点位于无穷远处。消失点也可能位于实际照片边界之外。
\end{checkedsolution}

\exercisequestion{SC4061-TUT07-Q03}{相机标定的主要步骤}

\textbf{对应讲义：}\relatedtopic{SC4061-L11-T04}{内外参与统一投影矩阵}；主讲义 pp.~26--32 的传统标定流程。

\textbf{题目。}描述 camera calibration（相机标定）的主要步骤。

\begin{checkedsolution}
标定利用已知的三维点及其二维像点估计相机参数。对世界坐标齐次向量 $\bm P_i^h=(X_i,Y_i,Z_i,1)^T$，针孔模型为
\[
  \lambda_i\begin{bmatrix}u_i\\v_i\\1\end{bmatrix}
    =\underbrace{K[R\mid t]}_{M\in\mathbb R^{3\times4}}\bm P_i^h.
\]
$K$ 为内参矩阵，$R,t$ 描述世界坐标到相机坐标的旋转和平移；$\lambda_i$ 是每个点各自的齐次比例因子。

\begin{enumerate}[label=\arabic*.]
  \item \textbf{准备已知几何的标定物。}使用精确制作的标定板或立方体，建立世界坐标系，确定角点的三维坐标并拍摄照片。
  \item \textbf{建立 3D--2D correspondences（三维到二维对应）。}用角点检测或边缘交点定位照片中的角点，记录像素坐标 $(u_i,v_i)$，与相应世界点逐一匹配。
  \item \textbf{估计投影矩阵 $M$。}消去每点的 $\lambda_i$ 后，一对点给出两个关于矩阵元素的线性约束；汇总全部对应点求解。讲义固定整体尺度，写成 $A\bm\theta=\bm b$；对应点较多时用伪逆求最小二乘解。
  \item \textbf{恢复内外参。}由估计的 $M$ 分解得到 $K,R,t$，完成讲义中的标定流程。
\end{enumerate}

\textbf{为何至少需要六对点？}一般 $3\times4$ 投影矩阵有 $12$ 个元素，整体非零倍数表示相同投影，因此有 $11$ 个独立自由度。每对点提供两个约束，数量下限为
\[
  2n\ge11\quad\Longrightarrow\quad\boxed{n\ge6}.
\]
达到数量下限还要求点的配置非退化，使方程提供足够独立的约束；通常使用更多、分布良好的对应点，以降低定位噪声的影响。

\textbf{适用条件。}单张照片中全部共面的点不足以确定一般的 $M$；平面标定板通常配合多个姿态和专门的平面标定方法。讲义令 $m_{34}=1$ 的写法要求原本 $m_{34}\ne0$，否则需换用其他尺度约束。镜头畸变参数属于实际标定的扩展，本题核心是上述针孔模型下的 $M,K,R,t$。
\end{checkedsolution}
\endgroup
